% !TEX program = xelatex
%-----------------------------------------------------------------------------------------------------------------------------------------------%
%	The MIT License (MIT)
%
%	Copyright (c) 2021 Jitin Nair
%
%	Permission is hereby granted, free of charge, to any person obtaining a copy
%	of this software and associated documentation files (the "Software"), to deal
%	in the Software without restriction, including without limitation the rights
%	to use, copy, modify, merge, publish, distribute, sublicense, and/or sell
%	copies of the Software, and to permit persons to whom the Software is
%	furnished to do so, subject to the following conditions:
%	
%	THE SOFTWARE IS PROVIDED "AS IS", WITHOUT WARRANTY OF ANY KIND, EXPRESS OR
%	IMPLIED, INCLUDING BUT NOT LIMITED TO THE WARRANTIES OF MERCHANTABILITY,
%	FITNESS FOR A PARTICULAR PURPOSE AND NONINFRINGEMENT. IN NO EVENT SHALL THE
%	AUTHORS OR COPYRIGHT HOLDERS BE LIABLE FOR ANY CLAIM, DAMAGES OR OTHER
%	LIABILITY, WHETHER IN AN ACTION OF CONTRACT, TORT OR OTHERWISE, ARISING FROM,
%	OUT OF OR IN CONNECTION WITH THE SOFTWARE OR THE USE OR OTHER DEALINGS IN
%	THE SOFTWARE.
%	
%
%-----------------------------------------------------------------------------------------------------------------------------------------------%

%----------------------------------------------------------------------------------------
%	DOCUMENT DEFINITION
%----------------------------------------------------------------------------------------

% article class because we want to fully customize the page and not use a cv template
\documentclass[a4paper,10pt]{article}

%----------------------------------------------------------------------------------------
%	FONT
%----------------------------------------------------------------------------------------

% % fontspec allows you to use TTF/OTF fonts directly
\usepackage{fontspec}
\defaultfontfeatures{Ligatures=TeX}

\setmainfont[
    Path = fonts/,
    UprightFont = *-Regular.ttf,
    BoldFont = *-Bold.ttf,
    ItalicFont = *-Italic.ttf,
    BoldItalicFont = *-BoldItalic.ttf
]{Poppins}


%----------------------------------------------------------------------------------------
%	PACKAGES
%----------------------------------------------------------------------------------------
\usepackage{url}
\usepackage{parskip} 	

%other packages for formatting
\RequirePackage{color}
\RequirePackage{graphicx}
\usepackage[usenames,dvipsnames]{xcolor}
\usepackage[
    top=1cm,       % Marge entre le haut de la feuille et votre nom
    bottom=1cm,    % Marge entre le bas de la feuille et vos intérêts
    left=1.1cm,      % Marge gauche
    right=1.1cm      % Marge droite
]{geometry}

%tabularx environment
\usepackage{tabularx}

%for lists within experience section
\usepackage{enumitem}

% centered version of 'X' col. type
\newcolumntype{C}{>{\centering\arraybackslash}X} 

%to prevent spillover of tabular into next pages
\usepackage{supertabular}
\usepackage{tabularx}
\newlength{\fullcollw}
\setlength{\fullcollw}{0.47\textwidth}

%custom \section
\usepackage{titlesec}				
\usepackage{multicol}
\usepackage{multirow}

%CV Sections inspired by: 
%http://stefano.italians.nl/archives/26
\titleformat{\section}{\Large\raggedright}{}{0em}{}[\titlerule]
\titlespacing{\section}{0pt}{5pt}{2pt}

%for publications
% \usepackage[style=authoryear,sorting=ynt, maxbibnames=2]{biblatex}

%Setup hyperref package, and colours for links
\usepackage[unicode, draft=false]{hyperref}
\definecolor{linkcolour}{rgb}{0,0.2,0.6}
\hypersetup{colorlinks,breaklinks,urlcolor=linkcolour,linkcolor=linkcolour}
% \addbibresource{citations.bib}
% \setlength\bibitemsep{1em}

%for social icons
\usepackage{fontawesome5}

% job listing environments
\newenvironment{jobshort}[2]
    {
    \begin{tabularx}{\linewidth}{@{}l X r@{}}
    \textbf{#1} & \hfill &  #2
    \end{tabularx}
    }
    {
    }

\newenvironment{joblong}[2]
    {
    \begin{tabularx}{\linewidth}{@{}l X r@{}}
    \textbf{#1} & \hfill &  #2
    \end{tabularx}
    \begin{minipage}[t]{\linewidth}
    \begin{itemize}[nosep,after=\strut, leftmargin=1em, itemsep=0pt,label=--]
    }
    {
    \end{itemize}
    \end{minipage}
    }

\newenvironment{project}[2]
    {
    \begin{tabularx}{\linewidth}{@{}l X r@{}}
    \textbf{#1} & \hfill & #2
    \end{tabularx}
    \begin{minipage}[t]{\linewidth}
    \begin{itemize}[nosep,after=\strut, leftmargin=1em, itemsep=0pt,label=--]
    }
    {
    \end{itemize}
    \end{minipage}
    }

%----------------------------------------------------------------------------------------
%	BEGIN DOCUMENT
%----------------------------------------------------------------------------------------
\begin{document}

% non-numbered pages
\pagestyle{empty} 

%----------------------------------------------------------------------------------------
%	TITLE
%----------------------------------------------------------------------------------------

\begin{tabularx}{\linewidth}{@{} C @{}}\\
\Huge{Maxence Debes} \\[12pt]
\href{https://maxencedbe.github.io}{\textcolor{black}{\raisebox{-0.05\height}\faGlobe\ maxencedbe.github.io}} \ | \
\href{https://linkedin.com/in/maxence-debes}{\textcolor{black}{\raisebox{-0.05\height}\faLinkedin\ maxence-debes}} \ | \
\href{https://github.com/maxencedbe}{\textcolor{black}{\raisebox{-0.05\height}\faGithub\ maxencedbe}} \\[0pt]
\href{mailto:maxence.debes@polytechnique.edu}{\textcolor{black}{\raisebox{-0.05\height}\faEnvelope\ maxence.debes@polytechnique.edu}} \ | \
\href{tel:+33781327922}{\textcolor{black}{\raisebox{-0.05\height}\faPhone\ +33\ 7\ 81\ 32\ 79\ 22}} \\[4pt]

\end{tabularx}

%Summary
\noindent
Étudiant en Master Data Science avec une solide formation en mathématiques et informatique. Spécialisé en machine learning avancé et en implémentation d'algorithmes d'IA de pointe.\\
Recherche d'un stage de fin d'études de 6 mois en IA/machine learning pour le printemps–été 2026.\\[-8pt]

%----------------------------------------------------------------------------------------
%	PROJECTS
%-----------------------------
\section{Projets}

\begin{project}{Hackathon Hi! Paris 2025 (2ème place)}
{\href{https://github.com/maxencedbe/hi_paris_2025}{\textcolor{black}{\raisebox{-0.05\height}{\large \faGithub}}}}
    \item Classé 2ème sur plus de 80 équipes lors de l'édition 2025 du hackathon de data science Hi!Paris. Objectif : prédire les scores PISA des élèves à partir de contextes socio-économiques complexes.
    \item Conception d'un modèle Gated (routeur composé de deux régresseurs XGBoost), atteignant un score $R^{2}$ de 0.79 sur le jeu de test final.
\end{project}\vspace{-6pt}

\begin{project}{Molecular Graph Captioning (compétition Kaggle)}
{\href{https://github.com/maxencedbe/molecular_graph_captioning}{\textcolor{black}{\raisebox{-0.05\height}{\large \faGithub}}}}
    \item Développement d'une architecture utilisant l'apprentissage contrastif pour aligner les embeddings de graphes moléculaires avec des descriptions en langage naturel dans un espace latent commun.
    \item Implémentation d'un encodeur de réseau de neurones de graphes dédié avec PyTorch Geometric, optimisé pour l'alignement sémantique via les métriques BERTScore et BLEU-4.
\end{project}\vspace{-6pt}

\begin{project}{DragonLLM --- Toxicité LLM \& Guardrails (en cours, Janvier 2026)}
{\href{https://github.com/vemz/capstone_dragonllm}{\textcolor{black}{\raisebox{-0.05\height}{\large \faGithub}}}}
    \item Objectif : entraîner un classifieur conçu pour interrompre la génération lorsque le modèle de langage commence à produire du contenu toxique ou nuisible.
    \item Implémentation des architectures Qwen3-Guardrails et inspirées des circuit-breakers pour l'application en temps réel de mesures de sécurité dans un modèle d'IA spécialisé pour le secteur financier.
\end{project}\vspace{-6pt}

\begin{project}{Protection des données \& anonymisation --- Projet Cassiopée}
{\href{https://github.com/maxencedbe/cassiopee}{\textcolor{black}{\raisebox{-0.05\height}{\large \faGithub}}}}
    \item Analyse des risques de ré-identification dans des données de santé anonymisées via OSINT, promouvant un usage éthique des données médicales pour la recherche en IA.
\end{project}\vspace{-6pt}

\begin{project}{Optimisation de parking automatisé}
{\href{https://github.com/maxencedbe/parking-optimization}{\textcolor{black}{\raisebox{-0.05\height}{\large \faGithub}}}}
    \item Conception d'un algorithme hybride combinant la recherche A* (pour la recherche de chemin) et le recuit simulé (pour l'optimisation globale) afin de minimiser les coûts de récupération.
\end{project} \\[-16pt]

%----------------------------------------------------------------------------------------
% EXPERIENCE SECTIONS
%----------------------------------------------------------------------------------------

%Experience
\section{Expérience professionnelle}

\begin{joblong}{Stagiaire Informatique --- Batigère Maison Familiale}{Juillet 2025 --- Août 2025}
\item Optimisation de l'infrastructure numérique par la migration de l'entreprise vers SharePoint afin d'améliorer la collaboration et la gestion documentaire, avec une adoption par 100\% des employés.
\item Conception et mise en œuvre d'une politique de conformité RGPD, renforçant la sécurité des données et assurant la conformité réglementaire.
\end{joblong}\vspace{-6pt}

\begin{joblong}{Stagiaire Développeur Java --- Atos}{Juillet 2024 --- Août 2024}
\item Conception et implémentation d'une application de suivi du télétravail grâce au générateur d'applications JHipster, facilitant l'accessibilité aux postes de travail.
\end{joblong} \\[-16pt]

%----------------------------------------------------------------------------------------
%	EDUCATION
%----------------------------------------------------------------------------------------
\section{Formation}
\begin{tabularx}{\linewidth}{@{}X r@{}}
\textbf{École Polytechnique --- Institut Polytechnique de Paris} & 2025 --- présent\\
\end{tabularx} \\
Master Data Science\\[1pt]
\textit{\textbf{Cours pertinents :} Advanced Learning for Text and Graph Data (ALTEGRAD --- Cours commun MVA), Deep Learning, Machine Learning, Reinforcement Learning, Representation Learning (Cours commun MVA), Optimisation pour la Data Science, Generative Models.} \\[6pt]
\begin{tabularx}{\linewidth}{@{}X r@{}}
\textbf{Télécom SudParis --- Institut Polytechnique de Paris} & 2023 --- présent\\
\end{tabularx}\\
Diplôme d'ingénieur en Mathématiques et Informatique\\[1pt]
\textit{\textbf{Cours pertinents :} Analyse réelle et Intégration, Statistiques et analyse de données, Optimisation, Calcul scientifique, Machine Learning et Data Mining, Statistiques appliquées, Processus stochastiques, Optimisation continue pour l'apprentissage.}\\[-8pt]

%----------------------------------------------------------------------------------------
%	SKILLS
%----------------------------------------------------------------------------------------
\section{Compétences \& Centres d'intérêt}

\textbf{Compétences techniques :} Python, PyTorch, Transformers, Scikit-learn, Pandas, SQL, Java, Git, Linux, \LaTeX\\
\textbf{Langues :} Français (natif), Anglais (C1 --- TOEIC 965/990), Allemand (A2)\\
\textbf{Engagement \& centres d'intérêt :} Ancien président d'association étudiante, handball (compétition), musique

\vfill
\end{document}